\section{Monte Carlo：从随机样本到数值答案}
\label{sec:05-Probability/08-Limit-Theorems/05}

你接到了一个任务。有人递来一个积分：

\[
I=\int_{0}^{1} e^{-x^3}\sin(x^2)\,dx.
\]

没有初等原函数。铺网格、梯形法则、Simpson——数值积分的老实工具可以上。把 \([0,1]\) 切成 \(m\) 段，每段用多项式逼近被积函数，误差大概 \(O(1/m^2)\) 或 \(O(1/m^4)\)（取决于方法的光滑性要求）。对一维平滑函数，这些方法快且准。

但换一个问题：十维空间里的积分。\(m^d\) 个网格点——\(m=100\)，\(d=10\)，那就是 \(100^{10}=10^{20}\) 个点。这个数字大到你再活一万年也算不完。

这不是数学家的胡思乱想。金融衍生品定价、Bayesian 后验期望、物理系统的配分函数、机器学习中的边缘似然——哪一个都在高维空间里求积分。网格方法的指数爆炸把路堵死了。

Monte Carlo 方法从另一扇门进来。

\subsection{重新定义问题：积分就是期望}

把积分改写一种形式：

\[
I=\int_{0}^{1}g(x)\cdot1\,dx
=\E[g(X)],\qquad X\sim\operatorname{Uniform}(0,1),
\]

其中 \(g(x)=e^{-x^3}\sin(x^2)\)。积分变成了一个随机变量 \(g(X)\) 的期望——\(X\) 在 \((0,1)\) 上均匀取值，被积函数 \(g\) 告诉你"这个 \(X\) 值对积分的贡献是多少"。

这个改写的妙处在于把"精确计算面积"替换成"估计一个平均值"。而估计平均值这件事，概率论有成熟的语言。

在 \((0,1)\) 上抽取 \(n\) 个独立均匀样本 \(X_1,\ldots,X_n\)，计算

\[
\widehat I_n=\frac1n\sum_{i=1}^{n}g(X_i).
\]

这就是 Monte Carlo 估计。\(n\) 个随机点，\(n\) 个函数值，一次算术平均——程序十行，算完收工。

直觉说这应该行得通：样本越多，平均越逼近期望。但直觉不告诉你在 \(n=1000\) 的时候离答案还有多远。你跑了一遍，拿到 \(\widehat I_{1000}=0.4327\)。你有多确定它是 0.43 而不是 0.40 或 0.434？

\subsection{大数定律说你会到——但不说你离多远}

第一个要确认的是：这个估计确实会收敛到目标值。大数定律回答的就是这一点。

若 \(\E[|g(X)|]<\infty\)，则

\[
\widehat I_n\to I
\]

依概率成立；在矩条件稍强（比如方差有限）时，几乎处处收敛也成立。Monte Carlo 有了理论基石：只要你能从目标分布中抽样，并且被积函数可积，估计值最终一定稳定在正确答案周围。

但这是定性的安慰。"最终会到"和"现在在哪儿"是两回事。你手握 \(\widehat I_{1000}=0.4327\)，大数定律既不会给你一个误差半径，也不会告诉你还需要加多少样本才能把误差压到某个阈值以下。

你把样本量翻到 \(n=4000\)，又跑了一次——\(\widehat I_{4000}=0.4331\)。两次相差 0.0004。这是偶然碰巧，还是典型波动？4000 个样本比 1000 个好多少？直觉说"好两倍左右"，但直觉也曾经说过地球是平的。

你需要一个定理，它不但说"你会收敛"，还给出收敛的快慢。

\subsection{中心极限定理给出误差的尺码}

CLT 告诉你误差的规模。设

\[
\sigma_g^2=\Var(g(X))<\infty,
\]

那么经过标准化后，估计误差服从渐近正态分布：

\[
\sqrt n\,(\widehat I_n-I)\xrightarrow{d}\mathcal N(0,\sigma_g^2).
\]

翻译成实用语言：当 \(n\) 不太小时（几百通常够用，具体看分布偏度），

\[
\widehat I_n-I\;\sim\;\mathcal N\!\left(0,\frac{\sigma_g^2}{n}\right).
\]

由此，标准误大约是

\[
\frac{\sigma_g}{\sqrt n}.
\]

这条公式里有两个因子。\(\sigma_g\) 是单次抽样 \(g(X_i)\) 的标准差——它衡量被积函数在你抽样分布下的自然波动。如果 \(g\) 在不同 \(x\) 处的取值差异巨大，\(\sigma_g\) 就大，需要更多样本才能把平均压稳。

\(1/\sqrt n\) 是你为增加样本付出的代价。要把标准误缩到原来的一半，样本量需要翻四倍；缩到原来的十分之一，样本量需要翻一百倍。\(\sqrt n\) 的减速比起 ODE 求解器的 \(O(h^4)\) 或数值积分的 \(O(1/m^2)\) 确实慢——但它有一个杀手级的优势：维数不出现在误差公式里。

在一维，20 个 Simpson 点可能比 10000 个 Monte Carlo 点还准。在十维，Simpson 需要至少 \(3^{10}\approx59000\) 个点才能搭起网格，Monte Carlo 还是只需要管那个 \(\sigma_g/\sqrt n\)。网格方法的收敛速度在维度升高的同时，底数也在变大——\(m^d\) 个点才能维持每个方向上同等密度。Monte Carlo 的 \(1/\sqrt n\) 和维数彻底脱钩。不是魔法，是概率不管你在多少维的空间里取平均。

\subsection{把你的运行结果变成一份附误差的答案}

\(\sigma_g\) 在实际中未知，用样本标准差替代：

\[
s_g^2=\frac1{n-1}\sum_{i=1}^{n}\bigl(g(X_i)-\widehat I_n\bigr)^2.
\]

标准误的估计值为

\[
\widehat{\mathrm{SE}}(\widehat I_n)\approx\frac{s_g}{\sqrt n}.
\]

配合正态近似，一个常用的 95\% 置信区间是

\[
\widehat I_n\pm1.96\,\frac{s_g}{\sqrt n}.
\]

现在你跑完 Monte Carlo 之后不再是一个光秃秃的数字。你可以说："积分约等于 \(0.433\pm0.008\)。"那半个宽度 \(0.008\) 来自 \(1.96\,s_g/\sqrt n\)。如果有人不满意这个精度，你可以反推需要多少样本：要 95\% 的半宽不超过 0.001，需要 \(n\ge(1.96\,s_g/0.001)^2\)，用当前 \(s_g\) 代进去，目标样本量就算出来了。

\subsection{用随机撒点估计 \(\pi\)——教科书级的示范}

在正方形 \([-1,1]^2\) 中均匀随机撒点 \((X_i,Y_i)\)。指示变量

\[
Z_i=\mathbf 1_{\{X_i^2+Y_i^2\le1\}}
\]

标记点是否落在内切单位圆内。落进去的概率是圆面积除以正方面积：

\[
P(Z_i=1)=\frac{\pi\cdot1^2}{2\cdot2}=\frac{\pi}{4}.
\]

因此

\[
\widehat\pi_n=4\bar Z_n=\frac4n\sum_{i=1}^{n}Z_i
\]

是 \(\pi\) 的无偏估计。每个 \(Z_i\) 是 \(\operatorname{Bernoulli}(\pi/4)\)，方差为

\[
\Var(\widehat\pi_n)=\frac{16\,p(1-p)}{n},\qquad p=\frac{\pi}{4}.
\]

取 \(n=10^6\)，大概拿到 3.141 级别的前三位。这不是计算 \(\pi\) 的高效方法——算术几何平均或 Ramanujan 级数每步多出好几位，秒杀 Monte Carlo。它的价值在另一个方向：把几何面积翻译成事件概率，然后用大数定律反向使用。\hyperref[sec:05-Probability/08-Limit-Theorems/02]{``大数定律与长期平均''}里用已知概率 \(p\) 预言长期频率；这里是反过来——用观测频率推未知概率。同一个联系，正反两个方向。

\subsection{方差缩减：比加样本更聪明的办法}

误差公式 \(\sigma_g/\sqrt n\) 里有两个可以下手的地方：要么增大 \(n\)，要么缩小 \(\sigma_g\)。加样本是直路，但如果你能把 \(\sigma_g\) 砍掉一半，等价于免费把样本量翻四倍。

想法和动机都简单——不要改变目标期望，换一种方式估计它，让新方式的波动更小。

\begin{itemize}
\tightlist
\item
  \textbf{控制变量}：找一个辅助函数 \(h(X)\)，期望 \(\E[h(X)]\) 已知（或者说你算得出来）。构造
  \[
  \widehat I_n^{\mathrm{CV}}=\frac1n\sum_{i=1}^{n}\bigl[g(X_i)-\beta(h(X_i)-\E[h(X)])\bigr].
  \]
  无论 \(\beta\) 取什么值，这个估计的期望仍然是 \(\E[g(X)]\)——你只是平均每个样本都减掉了零均值的东西。但如果 \(h\) 和 \(g\) 正相关，\(h\) 高于期望时 \(g\) 通常也偏高，减掉 \(\beta(h-\E[h])\) 等于主动消化掉这部分共同波动。最优的 \(\beta\) 是 \(\Cov(g,h)/\Var(h)\)，从样本中估计。

  举个例子：要估计 \(\E[e^X]\) 且 \(X\sim\operatorname{Uniform}(0,1)\)。\(h(X)=X\) 的期望是 \(1/2\)。\(e^X\) 和 \(X\) 正相关，把 \(X\) 对 \((0,1)\) 上的随机波动吸收掉一部分后，剩余部分的方差明显缩小。

\item
  \textbf{对偶变量}：让样本成对出现，每一对内部负相关。比如对 \((0,1)\) 上的均匀分布，抽 \(U\) 配对 \(1-U\)。\(g(U)\) 和 \(g(1-U)\) 如果一个随 \(U\) 递增、另一个随 \(1-U\) 递减，它们的平均值波动会小于两个独立抽样的平均。每个单独样本的边缘分布不变，但配对平均的过程里，峰谷互相填平。

\item
  \textbf{分层抽样}：把抽样空间切成若干子区域（层），每层内部独立抽样，再按层的概率权重合并。这保证了不会因为"运气不好"导致某个重要区域完全没有代表——每层都有至少几个样本。分层的精细程度和每层内部的样本分配策略共同决定了方差在层间和层内分别消除多少。

\item
  \textbf{重要性抽样}：原始问题是
  \[
  I=\int g(x)p(x)\,dx=\E_{X\sim p}[g(X)].
  \]
  引入另一个分布 \(q\)，改写为
  \[
  I=\int g(x)\frac{p(x)}{q(x)}q(x)\,dx
  =\E_{X\sim q}\!\left[g(X)\frac{p(X)}{q(X)}\right].
  \]
  从 \(q\) 抽样，用重要性权重 \(w(X)=p(X)/q(X)\) 修正偏差。如果 \(q\) 在 \(g(x)\) 对积分贡献大的地方更密集，那么每个样本的 \(g(X)w(X)\) 更均匀，方差更小。理想的 \(q^\star(x)\propto|g(x)|p(x)\) 让被积函数在 \(q\) 下近乎常数——但寻找 \(q^\star\) 通常比原始积分更难，实用的做法是构造一个可抽样且大致匹配的近似分布。

  重要性抽样的危险和它的功效一样大。如果 \(q\) 在某个区域密度极低而 \(g(x)p(x)\) 不小，少数几个样本带着暴涨的权重主导整个估计——方差可能变成无穷。选 \(q\) 时至少保证两件事：\(q(x)>0\) 覆盖所有 \(g(x)p(x)\neq0\) 的区域；尾部衰减不比 \(g(x)p(x)\) 慢。重要性权重过大的情况容易诊断——画直方图或记录最大权重比例。
\end{itemize}

这些技术全都在方差的来源上做文章。有的靠引入已知量的信息消化掉一部分随机性，有的靠重新分配抽样努力让每一份计算都投入在高信息区域。它们不替代大数定律和 CLT，而是在同样的极限定理框架下，让区间更窄。

\subsection{独立抽样是假设，不是现实}

前面全部推导假设 \(X_i\) 独立同分布。真实的计算环境里，这个假设会被两件事挑战。

\medskip
\noindent\textbf{伪随机数不是真随机。} 计算机产生的序列是确定性的，依靠统计检验来模拟独立均匀的样子。Mersenne Twister、PCG 等现代生成器在典型应用中足够好——周期长到你的计算在耗尽它之前会先耗尽你的寿命。但在极端的尾部概率或高维均匀性需求下，伪随机数的周期和格状结构可能留下系统偏差。这个问题在普通应用中可以忽略，但在需要"十万亿个样本"或"每个维度方向都均匀"的场景里要留个心。

\medskip
\noindent\textbf{MCMC 的相关性。} Bayesian 后验、高维分布的直接独立抽样常常不可行——退而求其次用 Markov 链模拟。链上的下一个样本依赖上一个，相邻样本正相关。\(n\) 个 MCMC 样本携带的信息少于 \(n\) 个独立样本。

定义一个量来量化这种损失——有效样本量。对平稳链上的某个函数值序列，近似公式为

\[
n_{\mathrm{eff}}\approx\frac{n}{1+2\sum_{k\ge1}\rho_k},
\]

其中 \(\rho_k\) 是该函数值在滞后 \(k\) 处的自相关。正自相关（MCMC 的常态）使 \(n_{\mathrm{eff}}<n\)；极端情况下链黏在某个区域迟迟不挪，几千步也只贡献几十个有效样本。

实际计算 \(n_{\mathrm{eff}}\) 需要截断自相关和（滞后太长时估计本身变得不可靠），常用的方法是找到自相关衰减到可忽略的窗口后截断。有效样本量依赖于你正在估计的那个具体函数——不同的函数在同一个链上可能有不同的自相关结构。

更根本的问题是：不要把"误差"和"Monte Carlo 误差"划等号。

\begin{itemize}
\tightlist
\item
  \textbf{Monte Carlo 随机误差}：可通过重复运行和标准误评估。这是你在控制的那部分。
\item
  \textbf{离散化偏差}：数值求解近似模型（如 Euler--Maruyama 离散化 SDE）引入的误差。它不是随机的——它是你选的近似方案本身偏离了真实连续过程。减小这个偏差要缩步长，与 Monte Carlo 误差互相独立。
\item
  \textbf{模型偏差}：目标分布本身不代表现实。再精确的积分算出来的也是模型里的量，不是真实现象的量。
\item
  \textbf{初始与收敛偏差}：MCMC 尚未进入平稳状态，链还在从初始点向典型集"热身"。前若干个样本携带的不是目标分布的信息，而是初始条件的信息。
\end{itemize}

样本均值非常稳定地收敛到错误的答案，仍然是错误的答案。这件事与\hyperref[sec:05-Probability/08-Limit-Theorems/02]{``大数定律与长期平均''}中谈到的选择偏差、共同系统误差同源：重复只能缩小随机波动，不能修正抽样机制本身的偏移。Monte Carlo 是一台精密机器——大数定律和 CLT 保证它在你的假设内部准确运转；如果假设本身就是歪的，精密的错误还是错误。

\subsection{一份可报告的 Monte Carlo 结果}

假设你负责把 Monte Carlo 结果写成报告。一份别人能够判断是否可信的报告至少包含：

\begin{enumerate}
\tightlist
\item
  目标期望的数学定义——在算什么；
\item
  样本从什么分布、通过什么机制产生（独立抽样、MCMC 加什么算法、或别的方案）；
\item
  样本量与独立性判断（若非独立，附有效样本量和自相关诊断）；
\item
  点估计；
\item
  标准误或附 CLT 依据的误差区间——你的不确定性自评；
\item
  是否使用了方差缩减技术，及其对估计量和对应方差估计量的冲击；
\item
  潜在系统偏差及其来源（离散化、模型错误设定、收敛不足等）。
\end{enumerate}

Monte Carlo 方法把概率论的核心定理从纸上的存在性声明搬到了桌面的可操作工具上。大数定律保证你终将到达，中心极限定理告诉你现在离目标大概还有多远。你不是在掷骰子碰运气——你在用定理度量不确定性，然后决定下一步是加样本、缩减方差、还是重新考虑你的抽样分布。
